% Standalone presentation wrapper; the retained proof body is unchanged.
\documentclass[11pt]{article}
\usepackage[a4paper,margin=25mm]{geometry}
\usepackage{amsmath,amssymb}
\usepackage[british]{babel}
\usepackage{url}
\usepackage[round,authoryear]{natbib}
\usepackage[hidelinks]{hyperref}
\setlength{\bibhang}{1.5em}
\setlength{\bibsep}{2pt}
\renewcommand{\bibfont}{\small}
\title{Supporting proofs for the bounded Chebyshev pricing estimator}
\author{Amanjeet Singh}
\date{9 October 2026}
\begin{document}
\maketitle
\noindent Keywords: quantum derivative pricing; Chebyshev sampling;
amplitude estimation; confidence coverage; logical resource bounds.

% BEGIN RETAINED CHEBYSHEV PROOF BODY
% Supporting fragment. The main paper supplies notation and references.
\section{A bounded reflection-only Chebyshev pricing estimator}
\label{supp:bounded-chebyshev-estimator}

I develop a finite, completely chargeable adaptation of the published
Chebyshev sampling framework of \citet{rallfuller2023}. The final published
ChebAE pseudocode permits repeated confidence looks at an unchanged degree.
For the adaptation below I spend failure probability on every fresh look,
including such repeats. I do not identify the bounded policy, its fixed
batch size or its fallback as the unmodified published algorithm. Nor do I
claim a new quantum signal-processing theorem.

\subsection{Exact logical signal and both parities}

Let $A$ be the already certified clean preparation of a finite financial
flag, with all private work initially zero. Write $P_g$ for its good-flag
projector, $p=\langle0|A^\dagger P_gA|0\rangle$, and
\[
 S_g=I-2P_g,\qquad S_0=I-2|0\rangle\langle0|,
 \qquad Q=-AS_0A^\dagger S_g.
\]
Here $|0\rangle$ refers to the invariant logical input, comparison and
flag register. Private workspace is excluded from this reflection only
because its clean-work invariance has already been proved. A measured
lookup architecture would additionally require its scalar-Kraus channel
proof, including the adjoint. The present comparison uses the retained
coherent unitary architecture.

For $0<p<1$ define the orthonormal vectors
$|g\rangle=P_gA|0\rangle/\sqrt p$ and
$|b\rangle=(I-P_g)A|0\rangle/\sqrt{1-p}$. Put
$\theta=\arcsin\sqrt p$ and
$|a\rangle=A|0\rangle=\sin\theta|g\rangle+
\cos\theta|b\rangle$. On the ordered basis $(|g\rangle,|b\rangle)$,
$S_g=\operatorname{diag}(-1,1)$, and
\[
 -AS_0A^\dagger=2|a\rangle\langle a|-I,
 \qquad
 Q=\begin{pmatrix}
 \cos(2\theta)&\sin(2\theta)\\
 -\sin(2\theta)&\cos(2\theta)
 \end{pmatrix}.
\]
The matrix follows by multiplying the displayed rank-one reflection by
$S_g$. Induction using the sine and cosine addition formulae proves
\[
 Q^k|a\rangle
 =\sin((2k+1)\theta)|g\rangle+
  \cos((2k+1)\theta)|b\rangle.
\]
This is the usual amplitude-amplification subspace calculation, rather
than a new primitive.

Let $T_d$ denote the degree-$d$ Chebyshev polynomial characterised by
$T_d(\cos x)=\cos(dx)$. For odd $d=2k+1$, apply $Q^kA$ and measure
the original good flag. Its success probability is
\[
 \sin^2(d\theta)
 =\cos^2(d(\pi/2-\theta))=T_d(\sqrt p)^2.
\]
For even $d=2k$, $k\geq1$, apply chronologically $A$, $Q^{k-1}$,
$S_g$, then $A^\dagger$, and observe the event that the invariant input
register is zero. The zero-event amplitude is
\begin{align*}
 \langle0|A^\dagger S_gQ^{k-1}A|0\rangle
 &=\langle a|S_gQ^{k-1}|a\rangle\\
 &=\cos\theta\cos((2k-1)\theta)
   -\sin\theta\sin((2k-1)\theta)\\
 &=\cos(2k\theta).
\end{align*}
Its squared modulus is $T_{2k}(\sqrt p)^2$. This construction uses
$d$ preparations or adjoints, $k$ good reflections and $k-1$ zero
reflections. The odd construction uses $d$ preparations or adjoints and
$k$ reflections of each kind. I retain these compositions in the charge;
the longer even circuit $A^\dagger Q^kA$ would require $d+2$ preparations
or adjoints. No saving from additional unevaluated fusion is presumed.

The uncontrolled global minus in $Q$ can be omitted from the physical
shot, because replacing $Q$ by $-Q$ contributes only $(-1)^k$ to the
whole final state. This omission is inapplicable to the controlled QPE
fallback, where the conditional minus is retained. All good and zero
reflection predicates are unchanged. For the even terminal event I
observe the logical input wires and compute their conjunction classically;
this has no additional coherent conjunction circuit. Every allocated
wire is nevertheless observed and conditionally reset before the next
fresh shot. At $p=0$ or $p=1$, the same probabilities follow directly by
applying the reflections to $A|0\rangle$; alternatively their values are
the continuous limits of the finite polynomial expressions. Thus no
singular normalisation of $|g\rangle$ or $|b\rangle$ is used at endpoints.

\subsection{A certified monotone inverse branch}

Set
\[
 t={2\over\pi}\arccos\sqrt p
 =1-{2\over\pi}\arcsin\sqrt p\in[0,1],
 \qquad q_d(p)=\cos^2(\pi dt/2).
\]
The map $p\mapsto t$ is decreasing. I enclose $t$ from the deterministic
two-sided probability prior using the existing directed rational sine
and inverse-sine calculation. Suppose the current outward enclosure is
$[t_-,t_+]$, and an integer $j\geq0$ satisfies
\[
 j\leq dt_-\leq dt_+\leq j+1.
 \tag{C1}\label{eq:cheb-monotone-branch}
\]
Write $dt=j+v$, $v\in[0,1]$, and let
$h(v)=\sin^2(\pi v/2)$. If $j$ is even then
$q_d=1-h(v)$; if $j$ is odd then $q_d=h(v)$. The derivative
$h'(v)=(\pi/2)\sin(\pi v)$ is positive for $0<v<1$, so each
restriction is one-to-one, including its endpoints. The zero derivative
at an endpoint does not invalidate this monotonicity argument.

If $[q_-,q_+]$ contains $q_d(p)$, let $h^{-1}(x)=
2\arcsin\sqrt x/\pi$. In an even branch the resulting exact $t$
interval is
\[
 \left[{j+h^{-1}(1-q_+)\over d},
       {j+h^{-1}(1-q_-)\over d}\right];
\]
in an odd branch it is
\[
 \left[{j+h^{-1}(q_-)\over d},
       {j+h^{-1}(q_+)\over d}\right].
\]
I replace each inverse endpoint by its separately certified outward
endpoint, then intersect with the preceding $t$ interval. Finally,
$p=1-h(t)$ gives the outward probability interval
\[
 [1-h(t_+)^{+},\,1-h(t_-)^{-}]
\]
intersected with the original deterministic prior. All intersections
retain containment whenever their input confidence events hold. An empty
intersection is recorded as an unsuccessful pilot, rather than repaired
by inserting an unproved point estimate.

The degree policy tests integers in descending order up to 2,047 and
accepts only a degree for which \eqref{eq:cheb-monotone-branch} is checked
by rational comparisons. A previous accepted degree remains valid on
every nonempty subset of its branch. After the first look the policy
changes degree only when the new accepted degree is at least twice the
previous one; otherwise it repeats the previous degree. This is a finite
implementation choice, not a proof of globally optimal degree selection.
Every terminal midpoint is formed from exact rational probability
endpoints. If their half-width is at most the required probability
tolerance $\varepsilon$, its error from the true $p$ is at most
$\varepsilon$ on the simultaneous confidence event.

\subsection{Fresh-batch confidence under adaptive decisions}

Let the total permitted failure probability be $\alpha\in(0,1)$.
There are at most twelve pilot looks. Each look uses exactly 100 fresh
complete shots, and spends $\alpha_{
\mathrm{look}}=\alpha/24$. A look reusing a degree is still a new look
for this allocation. I do not pool a random number of batches into a
fixed-sample interval. Conditional on the complete preceding history,
the selected degree is fixed and the new observations $X_1,\ldots,X_{100}$
are independent Bernoulli variables with common mean $q_d(p)$ under the
specified ideal logical experiment.

I use fixed-sample equal-tail binomial inversion. This is the established
Clopper--Pearson construction; I prove the relevant coverage here rather
than assuming that finite-sample coverage follows from a numerical root.
For $n=100$, $K=\sum_iX_i$ and $\gamma=\alpha_{\mathrm{look}}/2$, put
\[
 L_k(q)=\sum_{x=k}^n\binom nxq^x(1-q)^{n-x},\qquad
 U_k(q)=\sum_{x=0}^k\binom nxq^x(1-q)^{n-x}.
\]
For $k>0$, $L_k$ is strictly increasing on $(0,1)$ and has endpoint
values zero and one. One proof differentiates its finite sum and uses
$x\binom nx=n\binom{n-1}{x-1}$ and
$(n-x)\binom nx=n\binom{n-1}{x}$; cancellation gives
\[
 L_k'(q)=n\binom{n-1}{k-1}q^{k-1}(1-q)^{n-k}>0.
\]
For $k<n$, $U_k=1-L_{k+1}$ is strictly decreasing, with endpoint
values one and zero. Consequently the lower root $l_k$ of
$L_k(l_k)=\gamma$ exists uniquely for $k>0$, and the upper root
$u_k$ of $U_k(u_k)=\gamma$ exists uniquely for $k<n$. I set
$l_0=0$ and $u_n=1$.

Fix a true binomial parameter $q$. The sequence $L_k(q)$ decreases
with $k$. The event $q<l_K$ is therefore, when nonempty, the event
$K\geq k_*$ for the smallest $k_*$ such that $L_{k_*}(q)<\gamma$.
Its probability is $L_{k_*}(q)<\gamma$. If it is empty its probability
is zero. Similarly, $U_k(q)$ increases with $k$. The event $q>u_K$
is, when nonempty, $K\leq k^*$ for the largest $k^*$ such that
$U_{k^*}(q)<\gamma$, with probability below $\gamma$. Thus
\[
 \Pr_q\{q\notin[l_K,u_K]\}\leq2\gamma
 =\alpha_{\mathrm{look}}.
\]
The argument also treats $q=0,1$ directly, where the observed count is
deterministic and the interval contains its endpoint.

I isolate each root by exact dyadic bisection, retaining a bracket of
width at most $2^{-64}$. At an endpoint $q=z/2^b$, the numerator of a
binomial tail is the integer sum
\[
 \sum_{x\in I}\binom nxz^x(2^b-z)^{n-x},
\]
over the appropriate tail index set, with denominator $2^{bn}$.
Comparisons with rational $\gamma$ use integer cross multiplication,
so no floating-point root or unreported numerical tolerance enters.
The lower root is replaced by its bracket's lower endpoint and the
upper root by its bracket's upper endpoint. This enlarges the exact
confidence interval, preserving the proved coverage. Each interval is a
deterministic function of $k,n,\gamma$ and the fixed inverse precision;
the construction rule is registered before examining simulated outcomes.
Caching repeated counts changes only construction cost.

Conditional on the preceding history, this coverage argument applies
to the new fixed batch because its degree is already selected and its
parameter is $q_d(p)$. Integrating the conditional inequality over the preceding histories
bounds the unconditional failure probability of every reached look by
$\alpha/24$. The union of at most twelve such events therefore has
probability at most $\alpha/2$. Independence of different looks after
conditioning is not needed for that union bound. On its complement,
induction through the inverse-branch calculation proves containment of
the true $t$ and $p$ in every retained pilot interval. Adaptive stopping
therefore preserves the claimed probability error and confidence.

This statement requires fresh shots with the exact proved logical flag
and reflection channels. Retained simulated records or a finite
statistical audit do not demonstrate physical-device independence, gate
fault tolerance or empirical hardware execution. Those are separate
requirements for a physical result.

\subsection{Fallback coverage and a finite complete cost cap}

Before pilot execution I construct the canonical QPE fallback at the
same probability tolerance and the original deterministic prior, with
failure allowance $\alpha/2$. Let its complete retained logical T count
be $T_{\mathrm{fb}}$. The fallback includes its emitted Fourier tape,
confidence repetitions, controlled signs and reflections, observations
and resets. I set the pilot budget to
$B=\lfloor T_{\mathrm{fb}}/10\rfloor$.

Let $c_d$ be the completely charged T count for one shot of degree $d$
using the parity-specific construction above. Before each fresh batch
the controller checks that its accumulated count plus $100c_d$ is at
most $B$. If this check fails, the degree or look cap is reached, or a
confidence intersection is empty, I run the fixed QPE fallback with
fresh repetitions. The fallback answer is projected only onto the
original deterministic prior. I do not intersect its interval with a
pilot confidence interval, which could already have failed.

Let $E_P$ be the event that any pilot interval used in a terminal answer
fails, and let $E_F$ be failure of the fresh QPE fallback when reached.
The preceding result gives $\Pr(E_P)\leq\alpha/2$. Conditional on every
history reaching the fallback its fixed plan still has failure at most
$\alpha/2$, so $\Pr(E_F)\leq\alpha/2$. Every erroneous terminal answer
belongs to $E_P\cup E_F$. Hence the complete policy has probability
error at most $\varepsilon$ with confidence at least $1-\alpha$.
If the deterministic prior already has half-width at most
$\varepsilon$, its midpoint stops without quantum shots or confidence
expenditure.

The budget checks are independent of confidence validity and hold for
every possible history, including histories in which the intervals
fail. A successful pilot costs at most $B$. Every fallback history costs
at most
\[
 B+T_{\mathrm{fb}}\leq{11\over10}T_{\mathrm{fb}}.
\]
This is a finite worst-history guarantee, not an average over favourable
simulations. The comparator is the newly charged $\alpha/2$ fallback;
it is not the earlier $\alpha$ QPE cap, which can use fewer repetitions.
The bound concerns logical T gates. Classical control, preprocessing,
gate categories, observations, reset and memory retain their explicit
ledgers; I do not equate this inequality with physical elapsed time.

\subsection{An expected-gate break-even condition}

I fix the financial flag and its unknown true probability $p$. Let $S$
be the indicator that the pilot certifies its probability tolerance, and
let $C_P$ be its actual complete T charge. The randomness in this
paragraph concerns fresh shots at this fixed flag, rather than a
distribution over financial prices. If $T_0$ is the earlier full-$\alpha$
QPE charge, the complete hybrid charge has the pointwise identity
\[
 C=C_P+(1-S)T_{\mathrm{fb}}.
\]
All variables are bounded by the finite policy, so their expectations
exist. Linearity, without an independence assumption between $C_P$ and
$S$, gives
\[
 \mathbb EC=\mathbb EC_P+
       \{1-\Pr(S=1)\}T_{\mathrm{fb}}.
\]
Consequently the exact condition for a strict expected saving against
the earlier full-$\alpha$ comparator is
\begin{equation}
 \mathbb EC<T_0
 \quad\Longleftrightarrow\quad
 \Pr(S=1)>1-{T_0\over T_{\mathrm{fb}}}
                   +{\mathbb EC_P\over T_{\mathrm{fb}}}.
 \label{eq:cheb-expected-break-even}
\end{equation}
This remains valid when pilot success and pilot cost are strongly
dependent. It applies to any finitely charged pilot with a fixed
fallback, and follows from established elementary expectation rather
than a new quantum or statistical theorem.

In the four retained implementations the additional confidence
repetitions give the exact ratio $T_{\mathrm{fb}}/T_0=27/25$.
Since $0\leq C_P\leq B$, a necessary condition for a strict expected
saving is $\Pr(S=1)>2/27$. A sufficient condition is
\[
 \Pr(S=1)>{2\over27}+{B\over T_{\mathrm{fb}}}.
\]
The simpler uniform condition $\Pr(S=1)>47/270$ is also sufficient,
because $B/T_{\mathrm{fb}}\leq1/10$. The inequalities are strict:
equality in \eqref{eq:cheb-expected-break-even} gives equal expected
charges, not a saving. Equality at the conservative sufficient
threshold requires separately checking the actual mean pilot cost;
the budget envelope alone need not decide it. A success probability
above $2/27$ is not sufficient on its own. Even if the mean pilot cost
were zero, the necessary strict threshold would remain $2/27$.

The retained 400 synthetic histories use a fixed mesh of possible
probabilities and fixed pseudorandom seeds. Their success fractions and
pilot charges are descriptive properties of that mesh. They neither
identify $\Pr(S=1)$ at the unknown financial flag nor estimate
$\mathbb EC_P$ there. Thus they cannot establish that
\eqref{eq:cheb-expected-break-even} holds for the actual pricing task.
I retain the median-cost increases in all four diagnostic families;
neither an observed pilot success nor a smaller conditional path proves
an expected advantage.

For an undiscounted financial decoding satisfying
$|\widehat V-V|\leq\eta_{\mathrm{det}}+F\varepsilon$, I retain the
same deterministic allowance and exact monetary factor $F$ as the
existing R1 flag. Choosing $\varepsilon$ so that this sum is at most
the financial tolerance transfers the proved coverage to its derivative
price. No continuous classical price, sampled price or new financial
law enters this controller. Originality and algorithmic superiority
would require separate literature and complete performance comparisons;
the proved adaptation alone establishes neither.
% END RETAINED CHEBYSHEV PROOF BODY

\begin{thebibliography}{99}
\bibitem[Rall and Fuller(2023)]{rallfuller2023} Rall, P. and Fuller, B. (2023). Amplitude estimation from quantum signal processing. \emph{Quantum}, 7, 937. \url{https://doi.org/10.22331/q-2023-03-02-937}.
\end{thebibliography}
\end{document}
